\documentclass{ximera}
%% Macros shared by every activity. Each activity loads this file with
%% %% Macros shared by every activity. Each activity loads this file with
%% %% Macros shared by every activity. Each activity loads this file with
%% \input{../preamble.tex}, and the xourse file with \input{preamble.tex}.

%% MathML intent for screen readers: \xmunits, \xmabs, \xmpoint and the
%% other commands of TailorSwiftBot.sty. The file is latex/xmintent.sty.
\usepackage{xmintent}

%% The interactive pieces of Ximera: \answer, multipleChoice and \choice,
%% selectAll, hint and the rest. Keep this line: without it those commands
%% are undefined and come out as plain text. The file is
%% latex/ximera-components.sty.
\usepackage{ximera-components}

%% Tables with the tabularray package. The settings rowhead, xmcolhead,
%% caption and summary make their header cells known to screen readers, on
%% the web page and in the tagged PDF.
\usepackage{tabularray}

\newcommand{\RR}{\mathbb{R}}
\newcommand{\dfn}{\textit}
, and the xourse file with \input{preamble.tex}.

%% MathML intent for screen readers: \xmunits, \xmabs, \xmpoint and the
%% other commands of TailorSwiftBot.sty. The file is latex/xmintent.sty.
\usepackage{xmintent}

%% The interactive pieces of Ximera: \answer, multipleChoice and \choice,
%% selectAll, hint and the rest. Keep this line: without it those commands
%% are undefined and come out as plain text. The file is
%% latex/ximera-components.sty.
\usepackage{ximera-components}

%% Tables with the tabularray package. The settings rowhead, xmcolhead,
%% caption and summary make their header cells known to screen readers, on
%% the web page and in the tagged PDF.
\usepackage{tabularray}

\newcommand{\RR}{\mathbb{R}}
\newcommand{\dfn}{\textit}
, and the xourse file with \input{preamble.tex}.

%% MathML intent for screen readers: \xmunits, \xmabs, \xmpoint and the
%% other commands of TailorSwiftBot.sty. The file is latex/xmintent.sty.
\usepackage{xmintent}

%% The interactive pieces of Ximera: \answer, multipleChoice and \choice,
%% selectAll, hint and the rest. Keep this line: without it those commands
%% are undefined and come out as plain text. The file is
%% latex/ximera-components.sty.
\usepackage{ximera-components}

%% Tables with the tabularray package. The settings rowhead, xmcolhead,
%% caption and summary make their header cells known to screen readers, on
%% the web page and in the tagged PDF.
\usepackage{tabularray}

\newcommand{\RR}{\mathbb{R}}
\newcommand{\dfn}{\textit}


\title{First Steps}

\begin{document}
\begin{abstract}
One example of each basic kind of interactive question.
\end{abstract}
\maketitle

This page shows the building blocks. Copy any of them into your own activity.

Text can be \textbf{bold}, \emph{emphasised} or \underline{underlined}, and
underlining works in a formula too: $\underline{x+y}$.

\begin{problem}
An answer box inside a sentence: $2+2=\answer{4}$.
\end{problem}

\begin{problem}
An answer box in displayed math:
\[
\frac{1}{2}+\frac{1}{3}=\answer{\frac{5}{6}}.
\]
\end{problem}

\begin{problem}
A multiple choice question. Which of these numbers is prime?

\begin{multipleChoice}
\choice{$21$}
\choice[correct]{$23$}
\choice{$27$}
\end{multipleChoice}
\end{problem}

\begin{problem}
A question with more than one correct answer. Which of these numbers are
prime? Select all that apply.

\begin{selectAll}
\choice[correct]{$2$}
\choice{$9$}
\choice[correct]{$13$}
\choice{$15$}
\end{selectAll}
\end{problem}

\begin{problem}
A question with a hint. The set of real numbers is written $\RR$. How many
real solutions does $x^2=9$ have?

\begin{hint}
Remember that a negative number squared is positive.
\end{hint}

\[
\text{Number of solutions}=\answer{2}.
\]
\end{problem}

\begin{problem}
An answer box in a numbered equation:
\begin{equation}
3\cdot 4=\answer{12}
\end{equation}
\end{problem}

\begin{problem}
Answer boxes in aligned equations:
\begin{align*}
x+1&=5 \\
x&=\answer{4} \\
2x&=\answer{8}
\end{align*}
\end{problem}

\begin{problem}
An embedded Desmos graph, from the linear algebra textbook. Find a polynomial
function of an appropriate degree that passes through the points
\[
(-3,0),\quad (2,-2),\quad (-1,1).
\]
\[
f(x)=\answer{-0.3x^2-0.7x+0.6}
\]
Plot the graph of $f$ in the Desmos window below.

%% Not inside onlineOnly: that would leave the graph out of the PDF.
%% On its own, \desmos gives the graph on the web page and a link in the PDF.
\begin{center}
\desmos{aba3troiar}{950}{650}
\end{center}
\end{problem}

\end{document}
